\documentclass[12pt]{article}
\usepackage[utf8]{inputenc}
\usepackage[T1]{fontenc}
\usepackage{hyperref}
\title{Wumeng Yi Language Resources for K-12 Curriculum Development}
\author{Compiled by Claude Sonnet 5.5 (claude-sonnet-5-5) for Andr\'as Kornai}
\date{}
\begin{document}
\maketitle

\section{Introduction}
Wumeng Yi (ISO 639-3 ywm; Glottocode wume1236) is a Loloish language of the Sino-Tibetan family, spoken in northeast Yunnan, China. The family field in the request was garbled (``Bookkeeping''). The classification here comes from the search results, not from the request. Jamin Pelkey (1999), as cited in search results, counted about 38,300 Wumeng people in four counties of Zhaotong Prefecture. The language is described as part of the Eastern Yi group, related to Shuixi Nosu, Mangbu Nosu and Wusa Nasu, with limited mutual intelligibility. ``Wumeng'' is probably a name taken from the Wumeng Mountains. The searches found no dictionary or school materials specific to this language.

\section{Primary Educational Resources}
No primers, textbooks or school programmes specific to Wumeng were found.

\section{Linguistic and Reference Materials}
\begin{itemize}
\item Wikipedia, list of lesser-known Loloish languages: \url{https://en.wikipedia.org/wiki/List_of_lesser-known_Loloish_languages}
\item Wikipedia, Yi people: \url{https://en.wikipedia.org/wiki/Yi_people}
\item Wikipedia, Nisoish languages: \url{https://en.wikipedia.org/wiki/Nisoish_languages}
\item STEDT source record for the Yi--Chinese dictionary \emph{Yihan zidian} (1979). It covers Yi generally, not Wumeng specifically. \url{https://stedt.berkeley.edu/~stedt-cgi/rootcanal.pl/source/CSL-YIzd}
\item Unicode preliminary proposal on the Yi script: \url{https://www.unicode.org/L2/L2022/22183-yi-preliminary.pdf}
\end{itemize}

\section{Religious and Community Resources}
These are outsider profiles of the Wumeng people. They are not language resources.
\begin{itemize}
\item Joshua Project, Wumeng in China: \url{https://joshuaproject.net/index.php/people_groups/18713}
\item Asia Harvest profile, Wumeng: \url{https://people-groups.asiaharvest.org/people-group-profiles/yunnan/wumeng.pdf}
\end{itemize}

\section{Digital and Audio Resources}
No audio, video or software specific to Wumeng was found.

\section{Curriculum Development Context}
The searches found no evidence of Wumeng being taught in schools. Schooling in the area is in Mandarin. Yi-language teaching exists elsewhere, for example the Guizhou \emph{Concise Yi--Chinese Character Dictionary} revised in Bijie. Nothing found shows it covers Wumeng. A curriculum would need new documentation work with Wumeng speakers. Searches in Chinese (for example for Wumeng Yi dictionaries) were not carried out and might find more.

\end{document}